\chapter{Basic Usage}\label{svn.tour}

Theory is useful, but its application is just plain fun.
Let\textquotesingle s move now into the details of using Subversion. By
the time you reach the end of this chapter, you will be able to perform
all the tasks you need to use Subversion in a normal
day\textquotesingle s work. You\textquotesingle ll start with getting
your files into Subversion, followed by an initial checkout of your
code. We\textquotesingle ll then walk you through making changes and
examining those changes. You\textquotesingle ll also see how to bring
changes made by others into your working copy, examine them, and work
through any conflicts that might arise.

This chapter will not provide exhaustive coverage of all of
Subversion\textquotesingle s commands---rather, it\textquotesingle s a
conversational introduction to the most common Subversion tasks that
you\textquotesingle ll encounter. This chapter assumes that
you\textquotesingle ve read and understood
\hyperref[svn.basic]{Fundamental Concepts} and are familiar with the
general model of Subversion. For a complete reference of all commands,
see \hyperref[svn.ref.svn]{svn Reference---Subversion Command-Line
Client}.

Also, this chapter assumes that the reader is seeking information about
how to interact in a basic fashion with an existing Subversion
repository. No repository means no working copy; no working copy means
not much of interest in this chapter. There are many Internet sites
which offer free or inexpensive Subversion repository hosting services.
Or, if you\textquotesingle d prefer to set up and administer your own
repositories, check out \hyperref[svn.reposadmin]{Repository
Administration}. But don\textquotesingle t expect the examples in this
chapter to work without the user having access to a Subversion
repository.

Finally, any Subversion operation that contacts the repository over a
network may potentially require that the user authenticate. For the sake
of simplicity, our examples throughout this chapter avoid demonstrating
and discussing authentication. Be aware that if you hope to apply the
knowledge herein to an existing, real-world Subversion instance,
you\textquotesingle ll probably be forced to provide at least a username
and password to the server. See
\hyperref[svn.serverconfig.netmodel.creds]{Client Credentials} for a
detailed description of Subversion\textquotesingle s handling of
authentication and client credentials.

\section{Help!}\label{svn.tour.help}

It goes without saying that this book exists to be a source of
information and assistance for Subversion users new and old.
Conveniently, though, the Subversion command-line is self-documenting,
alleviating the need to grab a book off the shelf (wooden, virtual, or
otherwise). The \texttt{svn\ help} command is your gateway to that
built-in documentation:

\begin{verbatim}
$ svn help
usage: svn <subcommand> [options] [args]
Subversion command-line client, version 1.8.13.
Type 'svn help <subcommand>' for help on a specific subcommand.
Type 'svn --version' to see the program version and RA modules
  or 'svn --version --quiet' to see just the version number.

Most subcommands take file and/or directory arguments, recursing
on the directories.  If no arguments are supplied to such a
command, it recurses on the current directory (inclusive) by default.

Available subcommands:
   add
   blame (praise, annotate, ann)
   cat
…
\end{verbatim}

As described in the previous output, you can ask for help on a
particular subcommand by running \texttt{svn\ help\ SUBCOMMAND}.
Subversion will respond with the full usage message for that subcommand,
including its syntax, options, and behavior:

\begin{verbatim}
$ svn help help
help (?, h): Describe the usage of this program or its subcommands.
usage: help [SUBCOMMAND...]

Global options:
  --username ARG           : specify a username ARG
  --password ARG           : specify a password ARG
…
\end{verbatim}

Options and Switches and Flags, Oh My!

The Subversion command-line client has numerous command modifiers. Some
folks refer to such things as ``switches'' or ``flags''---in this book,
we\textquotesingle ll call them ``options''. You\textquotesingle ll find
the options supported by a given \texttt{svn} subcommand, plus a set of
options which are globally supported by all subcommands, listed near the
bottom of the built-in usage message for that subcommand.

Subversion\textquotesingle s options have two distinct forms: short
options are a single hyphen followed by a single letter, and long
options consist of two hyphens followed by several letters and hyphens
(e.g., \texttt{-s} and \texttt{-\/-this-is-a-long-option},
respectively). Every option has at least one long format. Some, such as
the \texttt{-\/-changelist} option, feature an abbreviated long-format
alias (\texttt{-\/-cl}, in this case). Only certain options---generally
the most-used ones---have an additional short format. To maintain
clarity in this book, we usually use the long form in code examples, but
when describing options, if there\textquotesingle s a short form,
we\textquotesingle ll provide the long form (to improve clarity) and the
short form (to make it easier to remember). Use the form
you\textquotesingle re more comfortable with when executing your own
Subversion commands.

Many Unix-based distributions of Subversion include manual pages of the
sort that can be invoked using the \texttt{man} program, but those tend
to carry only pointers to other sources of real help, such as the
project\textquotesingle s website and to the website which hosts this
book. Also, several companies offer Subversion help and support, too,
usually via a mixture of web-based discussion forums and fee-based
consulting. And of course, the Internet holds a decade\textquotesingle s
worth of Subversion-related discussions just begging to be located by
your favorite search engine. Subversion help is never too far away.

\section{Getting Data into Your Repository}\label{svn.tour.importing}

You can get new files into your Subversion repository in two ways:
\texttt{svn\ import} and \texttt{svn\ add}. We\textquotesingle ll
discuss \texttt{svn\ import} now and will discuss \texttt{svn\ add}
later in this chapter when we review a typical day with Subversion.

\subsection{Importing Files and
Directories}\label{svn.tour.importing.import}

The \texttt{svn\ import} command is a quick way to copy an unversioned
tree of files into a repository, creating intermediate directories as
necessary. \texttt{svn\ import} doesn\textquotesingle t require a
working copy, and your files are immediately committed to the
repository. You typically use this when you have an existing tree of
files that you want to begin tracking in your Subversion repository. For
example:

\begin{verbatim}
$ svn import /path/to/mytree \
             http://svn.example.com/svn/repo/some/project \
             -m "Initial import"
Adding         mytree/foo.c
Adding         mytree/bar.c
Adding         mytree/subdir
Adding         mytree/subdir/quux.h

Committed revision 1.
$
\end{verbatim}

The previous example copied the contents of the local directory
\texttt{mytree} into the directory \texttt{some/project} in the
repository. Note that you didn\textquotesingle t have to create that new
directory first---\texttt{svn\ import} does that for you. Immediately
after the commit, you can see your data in the repository:

\begin{verbatim}
$ svn list http://svn.example.com/svn/repo/some/project
bar.c
foo.c
subdir/
$
\end{verbatim}

Note that after the import is finished, the original local directory is
\emph{not} converted into a working copy. To begin working on that data
in a versioned fashion, you still need to create a fresh working copy of
that tree.

\subsection{Recommended Repository
Layout}\label{svn.tour.importing.layout}

Subversion provides the ultimate flexibility in terms of how you arrange
your data. Because it simply versions directories and files, and because
it ascribes no particular meaning to any of those objects, you may
arrange the data in your repository in any way that you choose.
Unfortunately, this flexibility also means that it\textquotesingle s
easy to find yourself ``lost without a roadmap'' as you attempt to
navigate different Subversion repositories which may carry completely
different and unpredictable arrangements of the data within them.

To counteract this confusion, we recommend that you follow a repository
layout convention (established long ago, in the nascency of the
Subversion project itself) in which a handful of strategically named
Subversion repository directories convey valuable meaning about the data
they hold. Most projects have a recognizable ``main line'', or trunk, of
development; some branches, which are divergent copies of development
lines; and some tags, which are named, stable snapshots of a particular
line of development. So we first recommend that each project have a
recognizable project root in the repository, a directory under which all
of the versioned information for that project---and only that
project---lives. Secondly, we suggest that each project root contain a
\texttt{trunk} subdirectory for the main development line, a
\texttt{branches} subdirectory in which specific branches (or
collections of branches) will be created, and a \texttt{tags}
subdirectory in which specific tags (or collections of tags) will be
created. Of course, if a repository houses only a single project, the
root of the repository can serve as the project root, too.

Here are some examples:

\begin{verbatim}
$ svn list file:///var/svn/single-project-repo
trunk/
branches/
tags/
$ svn list file:///var/svn/multi-project-repo
project-A/
project-B/
$ svn list file:///var/svn/multi-project-repo/project-A
trunk/
branches/
tags/
$
\end{verbatim}

We talk much more about tags and branches in
\hyperref[svn.branchmerge]{Branching and Merging}. For details and some
advice on how to set up repositories when you have multiple projects,
see \hyperref[svn.branchmerge.maint.layout]{Repository Layout}. Finally,
we discuss project roots more in
\hyperref[svn.reposadmin.projects.chooselayout]{Planning Your Repository
Organization}.

\subsection{What\textquotesingle s In a
Name?}\label{svn.tour.importing.naming}

Subversion tries hard not to limit the type of data you can place under
version control. The contents of files and property values are stored
and transmitted as binary data, and
\hyperref[svn.advanced.props.special.mime-type]{File Content Type} tells
you how to give Subversion a hint that ``textual'' operations
don\textquotesingle t make sense for a particular file. There are a few
places, however, where Subversion places restrictions on information it
stores.

Subversion internally handles certain bits of data---for example,
property names, pathnames, and log messages---as UTF-8-encoded Unicode.
This is not to say that all your interactions with Subversion must
involve UTF-8, though. As a general rule, Subversion clients will
gracefully and transparently handle conversions between UTF-8 and the
encoding system in use on your computer, if such a conversion can
meaningfully be done (which is the case for most common encodings in use
today).

In WebDAV exchanges and older versions of some of
Subversion\textquotesingle s administrative files, paths are used as XML
attribute values, and property names in XML tag names. This means that
pathnames can contain only legal XML (1.0) characters, and properties
are further limited to ASCII characters. Subversion also prohibits
\texttt{TAB}, \texttt{CR}, and \texttt{LF} characters in path names to
prevent paths from being broken up in diffs or in the output of commands
such as \texttt{svn\ log} or \texttt{svn\ status}.

While it may seem like a lot to remember, in practice these limitations
are rarely a problem. As long as your locale settings are compatible
with UTF-8 and you don\textquotesingle t use control characters in path
names, you should have no trouble communicating with Subversion. The
command-line client adds an extra bit of help---to create ``legally
correct'' versions for internal use it will automatically escape illegal
path characters as needed in URLs that you type.

Of course, when it comes to choosing valid path names, Subversion
isn\textquotesingle t the only limiting factor. Teams using multiple
operating systems need to consider the limitations placed on path names
by those operating systems, too. For example, while Windows disallows
the use of colon characters in file names, a user on a Linux system can
very easily add such a file to version control, resulting in a dataset
that can no longer be checked out on Windows. Adding multiple files to a
directory whose names differ only in their letter casing will likewise
cause problems for users checking out working copies onto
case-insensitive filesystems. So, some broad awareness of the various
limitations introduced by different operating systems and filesystems,
then, is recommended.

\section{Creating a Working Copy}\label{svn.tour.initial}

Most of the time, you will start using a Subversion repository by
performing a checkout of your project. Checking out a directory from a
repository creates a working copy of that directory on your local
machine. Unless otherwise specified, this copy contains the youngest
(that is, most recently created or modified) versions of the directory
and its children found in the Subversion repository:

\begin{verbatim}
$ svn checkout http://svn.example.com/svn/repo/trunk
A    trunk/README
A    trunk/INSTALL
A    trunk/src/main.c
A    trunk/src/header.h
…
Checked out revision 8810.
$
\end{verbatim}

Although the preceding example checks out the trunk directory, you can
just as easily check out a deeper subdirectory of a repository by
specifying that subdirectory\textquotesingle s URL as the checkout URL:

\begin{verbatim}
$ svn checkout http://svn.example.com/svn/repo/trunk/src
A    src/main.c
A    src/header.h
A    src/lib/helpers.c
…
Checked out revision 8810.
$
\end{verbatim}

Since Subversion uses a copy-modify-merge model instead of
lock-modify-unlock (see \hyperref[svn.basic.vsn-models]{Versioning
Models}), you can immediately make changes to the files and directories
in your working copy. Your working copy is just like any other
collection of files and directories on your system. You can edit the
files inside it, rename it, even delete the entire working copy and
forget about it.

While your working copy is ``just like any other collection of files and
directories on your system,'' you can edit files at will, but you must
tell Subversion about \emph{everything else} that you do. For example,
if you want to copy or move an item in a working copy, you should use
\texttt{svn\ copy} or \texttt{svn\ move} instead of the copy and move
commands provided by your operating system. We\textquotesingle ll talk
more about them later in this chapter.

Unless you\textquotesingle re ready to commit the addition of a new file
or directory or changes to existing ones, there\textquotesingle s no
need to further notify the Subversion server that you\textquotesingle ve
done anything.

What Is This .svn Directory?

The topmost directory of a working copy---and prior to version 1.7,
every versioned subdirectory thereof---contains a special administrative
subdirectory named \texttt{.svn}. Usually, your operating
system\textquotesingle s directory listing commands
won\textquotesingle t show this subdirectory, but it is nevertheless an
important directory. Whatever you do, don\textquotesingle t delete or
change anything in the administrative area! Subversion uses that
directory and its contents to manage your working copy.

Notice that in the previous pair of examples, Subversion chose to create
a working copy in a directory named for the final component of the
checkout URL. This occurs only as a convenience to the user when the
checkout URL is the only bit of information provided to the
\texttt{svn\ checkout} command. Subversion\textquotesingle s
command-line client gives you additional flexibility, though, allowing
you to optionally specify the local directory name that Subversion
should use for the working copy it creates. For example:

\begin{verbatim}
$ svn checkout http://svn.example.com/svn/repo/trunk my-working-copy
A    my-working-copy/README
A    my-working-copy/INSTALL
A    my-working-copy/src/main.c
A    my-working-copy/src/header.h
…
Checked out revision 8810.
$
\end{verbatim}

If the local directory you specify doesn\textquotesingle t yet exist,
that\textquotesingle s okay---\texttt{svn\ checkout} will create it for
you.

\section{Basic Work Cycle}\label{svn.tour.cycle}

Subversion has numerous features, options, bells, and whistles, but on a
day-to-day basis, odds are that you will use only a few of them. In this
section, we\textquotesingle ll run through the most common things that
you might find yourself doing with Subversion in the course of a
day\textquotesingle s work.

The typical work cycle looks like this:

\begin{enumerate}
\def\labelenumi{\arabic{enumi}.}
\item
  \emph{Update your working copy.} This involves the use of the
  \texttt{svn\ update} command.
\item
  \emph{Make your changes.} The most common changes that
  you\textquotesingle ll make are edits to the contents of your existing
  files. But sometimes you need to add, remove, copy and move files and
  directories---the \texttt{svn\ add}, \texttt{svn\ delete},
  \texttt{svn\ copy}, and \texttt{svn\ move} commands handle those sorts
  of structural changes within the working copy.
\item
  \emph{Review your changes.} The \texttt{svn\ status} and
  \texttt{svn\ diff} commands are critical to reviewing the changes
  you\textquotesingle ve made in your working copy.
\item
  \emph{Fix your mistakes.} Nobody\textquotesingle s perfect, so as you
  review your changes, you may spot something that\textquotesingle s not
  quite right. Sometimes the easiest way to fix a mistake is start all
  over again from scratch. The \texttt{svn\ revert} command restores a
  file or directory to its unmodified state.
\item
  \emph{Resolve any conflicts (merge others\textquotesingle{} changes).}
  In the time it takes you to make and review your changes, others might
  have made and published changes, too. You\textquotesingle ll want to
  integrate their changes into your working copy to avoid the potential
  out-of-dateness scenarios when you attempt to publish your own. Again,
  the \texttt{svn\ update} command is the way to do this. If this
  results in local conflicts, you\textquotesingle ll need to resolve
  those using the \texttt{svn\ resolve} command.
\item
  \emph{Publish (commit) your changes.} The \texttt{svn\ commit} command
  transmits your changes to the repository where, if they are accepted,
  they create the newest versions of all the things you modified. Now
  others can see your work, too!
\end{enumerate}

\subsection{Update Your Working Copy}\label{svn.tour.cycle.update}

When working on a project that is being modified via multiple working
copies, you\textquotesingle ll want to update your working copy to
receive any changes committed from other working copies since your last
update. These might be changes that other members of your project team
have made, or they might simply be changes you\textquotesingle ve made
yourself from a different computer. To protect your data, Subversion
won\textquotesingle t allow you commit new changes to out-of-date files
and directories, so it\textquotesingle s best to have the latest
versions of all your project\textquotesingle s files and directories
before making new changes of your own.

Use \texttt{svn\ update} to bring your working copy into sync with the
latest revision in the repository:

\begin{verbatim}
$ svn update
Updating '.':
U    foo.c
U    bar.c
Updated to revision 2.
$
\end{verbatim}

In this case, it appears that someone committed modifications to both
\texttt{foo.c} and \texttt{bar.c} since the last time you updated, and
Subversion has updated your working copy to include those changes.

When the server sends changes to your working copy via
\texttt{svn\ update}, a letter code is displayed next to each item to
let you know what actions Subversion performed to bring your working
copy up to date. To find out what these letters mean, run
\texttt{svn\ help\ update} or see \hyperref[svn.ref.svn.c.update]{svn
update (up)} in \hyperref[svn.ref.svn]{svn Reference---Subversion
Command-Line Client}.

\subsection{Make Your Changes}\label{svn.tour.cycle.edit}

Now you can get to work and make changes in your working copy. You can
make two kinds of changes to your working copy: file changes and tree
changes. You don\textquotesingle t need to tell Subversion that you
intend to change a file; just make your changes using your text editor,
word processor, graphics program, or whatever tool you would normally
use. Subversion automatically detects which files have been changed, and
in addition, it handles binary files just as easily as it handles text
files---and just as efficiently, too. Tree changes are different, and
involve changes to a directory\textquotesingle s structure. Such changes
include adding and removing files, renaming files or directories, and
copying files or directories to new locations. For tree changes, you use
Subversion operations to ``schedule'' files and directories for removal,
addition, copying, or moving. These changes may take place immediately
in your working copy, but no additions or removals will happen in the
repository until you commit them.

Versioning Symbolic Links

On non-Windows platforms, Subversion is able to version files of the
special type symbolic link (or ``symlink''). A symlink is a file that
acts as a sort of transparent reference to some other object in the
filesystem, allowing programs to read and write to those objects
indirectly by performing operations on the symlink itself.

When a symlink is committed into a Subversion repository, Subversion
remembers that the file was in fact a symlink, as well as the object to
which the symlink ``points.'' When that symlink is checked out to
another working copy on a non-Windows system, Subversion reconstructs a
real filesystem-level symbolic link from the versioned symlink. But that
doesn\textquotesingle t in any way limit the usability of working copies
on systems such as Windows that do not support symlinks. On such
systems, Subversion simply creates a regular text file whose contents
are the path to which the original symlink pointed. While that file
can\textquotesingle t be used as a symlink on a Windows system, it also
won\textquotesingle t prevent Windows users from performing their other
Subversion-related activities.

Here is an overview of the five Subversion subcommands that
you\textquotesingle ll use most often to make tree changes:

\begin{description}
\item[\texttt{svn\ add\ FOO}]
Use this to schedule the file, directory, or symbolic link \texttt{FOO}
to be added to the repository. When you next commit, \texttt{FOO} will
become a child of its parent directory. Note that if \texttt{FOO} is a
directory, everything underneath \texttt{FOO} will be scheduled for
addition. If you want only to add \texttt{FOO} itself, pass the
\texttt{-\/-depth=empty} option.
\item[\texttt{svn\ delete\ FOO}]
Use this to schedule the file, directory, or symbolic link \texttt{FOO}
to be deleted from the repository. \texttt{FOO} is immediately deleted
from your working copy. (Of course, nothing is ever totally deleted from
the repository---just from its \texttt{HEAD} revision. You may continue
to access the deleted item in previous revisions). \footnote{Should you
  desire to resurrect the item so that it is again present in
  \texttt{HEAD}, see
  \hyperref[svn.branchmerge.basicmerging.resurrect]{Resurrecting Deleted
  Items}.}
\item[\texttt{svn\ copy\ FOO\ BAR}]
Create a new item \texttt{BAR} as a duplicate of \texttt{FOO} and
automatically schedule \texttt{BAR} for addition. When \texttt{BAR} is
added to the repository on the next commit, its copy history is recorded
(as having originally come from \texttt{FOO}). \texttt{svn\ copy} does
not create intermediate directories unless you pass the
\texttt{-\/-parents} option.
\item[\texttt{svn\ move\ FOO\ BAR}]
This command is exactly the same as running
\texttt{svn\ copy\ FOO\ BAR;\ svn\ delete\ FOO}. That is, \texttt{BAR}
is scheduled for addition as a copy of \texttt{FOO}, and \texttt{FOO} is
scheduled for removal. \texttt{svn\ move} does not create intermediate
directories unless you pass the \texttt{-\/-parents} option.
\item[\texttt{svn\ mkdir\ FOO}]
This command is exactly the same as running
\texttt{mkdir\ FOO;\ svn\ add\ FOO}. That is, a new directory named
\texttt{FOO} is created and scheduled for addition.
\end{description}

Changing the Repository Without a Working Copy

Subversion \emph{does} offer ways to immediately commit tree changes to
the repository without an explicit commit action. In particular,
specific uses of \texttt{svn\ mkdir}, \texttt{svn\ copy},
\texttt{svn\ move}, and \texttt{svn\ delete} can operate directly on
repository URLs as well as on working copy paths. Of course, as
previously mentioned, \texttt{svn\ import} always makes direct changes
to the repository. We discuss the ways to commit tree changes without a
working copy in \hyperref[svn.advanced.working-without-a-wc]{Working
Without a Working Copy}.

There are pros and cons to performing URL-based operations. One obvious
advantage to doing so is speed: sometimes, checking out a working copy
that you don\textquotesingle t already have solely to perform some
seemingly simple action is an overbearing cost. A disadvantage is that
you are generally limited to a single, or single type of, operation at a
time when operating directly on URLs. Finally, the primary advantage of
a working copy is in its utility as a sort of ``staging area'' for
changes. You can make sure that the changes you are about to commit make
sense in the larger scope of your project before committing them. And,
of course, these staged changes can be as complex or as a simple as they
need to be, yet result in but a single new revision when committed.

\subsection{Review Your Changes}\label{svn.tour.cycle.examine}

Once you\textquotesingle ve finished making changes, you need to commit
them to the repository, but before you do so, it\textquotesingle s
usually a good idea to take a look at exactly what
you\textquotesingle ve changed. By examining your changes before you
commit, you can compose a more accurate log message (a human-readable
description of the committed changes stored alongside those changes in
the repository). You may also discover that you\textquotesingle ve
inadvertently changed a file, and that you need to undo that change
before committing. Additionally, this is a good opportunity to review
and scrutinize changes before publishing them. You can see an overview
of the changes you\textquotesingle ve made by using the
\texttt{svn\ status} command, and you can dig into the details of those
changes by using the \texttt{svn\ diff} command.

Look Ma! No Network!

You can use the commands \texttt{svn\ status}, \texttt{svn\ diff}, and
\texttt{svn\ revert} without any network access even if your repository
\emph{is} across the network. This makes it easy to manage and review
your changes-in-progress when you are working offline or are otherwise
unable to contact your repository over the network.

Subversion does this by keeping private caches of pristine, unmodified
versions of each versioned file inside its working copy administrative
area (or prior to version 1.7, potentially multiple administrative
areas). This allows Subversion to report---and revert---local
modifications to those files \emph{without network access}. This cache
(called the text-base) also allows Subversion to send the
user\textquotesingle s local modifications during a commit to the server
as a compressed delta (or ``difference'') against the pristine version.
Having this cache is a tremendous benefit---even if you have a fast
Internet connection, it\textquotesingle s generally much faster to send
only a file\textquotesingle s changes rather than the whole file to the
server.

\subsubsection{See an overview of your
changes}\label{svn.tour.cycle.examine.status}

To get an overview of your changes, use the \texttt{svn\ status}
command. You\textquotesingle ll probably use \texttt{svn\ status} more
than any other Subversion command.

Because the \texttt{cvs\ status} command\textquotesingle s output was so
noisy, and because \texttt{cvs\ update} not only performs an update, but
also reports the status of your local changes, most CVS users have grown
accustomed to using \texttt{cvs\ update} to report their changes. In
Subversion, the update and status reporting facilities are completely
separate.

If you run \texttt{svn\ status} at the top of your working copy with no
additional arguments, it will detect and report all file and tree
changes you\textquotesingle ve made.

\begin{verbatim}
$ svn status
?       scratch.c
A       stuff/loot
A       stuff/loot/new.c
D       stuff/old.c
M       bar.c
$
\end{verbatim}

In its default output mode, \texttt{svn\ status} prints seven columns of
characters, followed by several whitespace characters, followed by a
file or directory name. The first column tells the status of a file or
directory and/or its contents. Some of the most common codes that
\texttt{svn\ status} displays are:

\begin{description}
\item[\texttt{?\ item}]
The file, directory, or symbolic link \texttt{item} is not under version
control.
\item[\texttt{A\ item}]
The file, directory, or symbolic link \texttt{item} has been scheduled
for addition into the repository.
\item[\texttt{C\ item}]
The file \texttt{item} is in a state of conflict. That is, changes
received from the server during an update overlap with local changes
that you have in your working copy (and weren\textquotesingle t resolved
during the update). You must resolve this conflict before committing
your changes to the repository.
\item[\texttt{D\ item}]
The file, directory, or symbolic link \texttt{item} has been scheduled
for deletion from the repository.
\item[\texttt{M\ item}]
The contents of the file \texttt{item} have been modified.
\end{description}

If you pass a specific path to \texttt{svn\ status}, you get information
about that item alone:

\begin{verbatim}
$ svn status stuff/fish.c
D       stuff/fish.c
\end{verbatim}

\texttt{svn\ status} also has a \texttt{-\/-verbose} (\texttt{-v})
option, which will show you the status of \emph{every} item in your
working copy, even if it has not been changed:

\begin{verbatim}
$ svn status -v
M               44        23    sally     README
                44        30    sally     INSTALL
M               44        20    harry     bar.c
                44        18    ira       stuff
                44        35    harry     stuff/trout.c
D               44        19    ira       stuff/fish.c
                44        21    sally     stuff/things
A                0         ?     ?        stuff/things/bloo.h
                44        36    harry     stuff/things/gloo.c
\end{verbatim}

This is the ``long form'' output of \texttt{svn\ status}. The letters in
the first column mean the same as before, but the second column shows
the working revision of the item. The third and fourth columns show the
revision in which the item last changed, and who changed it.

None of the prior invocations to \texttt{svn\ status} contact the
repository---they merely report what is known about the working copy
items based on the records stored in the working copy administrative
area and on the timestamps and contents of modified files. But sometimes
it is useful to see which of the items in your working copy have been
modified in the repository since the last time you updated your working
copy. For this, \texttt{svn\ status} offers the
\texttt{-\/-show-updates} (\texttt{-u}) option, which contacts the
repository and adds information about items that are out of date:

\begin{verbatim}
$ svn status -u -v
M      *        44        23    sally     README
M               44        20    harry     bar.c
       *        44        35    harry     stuff/trout.c
D               44        19    ira       stuff/fish.c
A                0         ?     ?        stuff/things/bloo.h
Status against revision:   46
\end{verbatim}

Notice in the previous example the two asterisks: if you were to run
\texttt{svn\ update} at this point, you would receive changes to
\texttt{README} and \texttt{trout.c}. This tells you some very useful
information---because one of those items is also one that you have
locally modified (the file \texttt{README}), you\textquotesingle ll need
to update and get the server\textquotesingle s changes for that file
before you commit, or the repository will reject your commit for being
out of date. We discuss this in more detail later.

\texttt{svn\ status} can display much more information about the files
and directories in your working copy than we\textquotesingle ve shown
here---for an exhaustive description of \texttt{svn\ status} and its
output, run \texttt{svn\ help\ status} or see
\hyperref[svn.ref.svn.c.status]{svn status (stat, st)} in
\hyperref[svn.ref.svn]{svn Reference---Subversion Command-Line Client}.

\subsubsection{Examine the details of your local
modifications}\label{svn.tour.cycle.examine.diff}

Another way to examine your changes is with the \texttt{svn\ diff}
command, which displays differences in file content. When you run
\texttt{svn\ diff} at the top of your working copy with no arguments,
Subversion will print the changes you\textquotesingle ve made to
human-readable files in your working copy. It displays those changes in
unified diff format, a format which describes changes as ``hunks'' (or
``snippets'') of a file\textquotesingle s content where each line of
text is prefixed with a single-character code: a space, which means the
line was unchanged; a minus sign (\texttt{-}), which means the line was
removed from the file; or a plus sign (\texttt{+}), which means the line
was added to the file. In the context of \texttt{svn\ diff}, those
minus-sign- and plus-sign-prefixed lines show how the lines looked
before and after your modifications, respectively.

Here\textquotesingle s an example:

\begin{verbatim}
$ svn diff
Index: bar.c
===================================================================
--- bar.c   (revision 3)
+++ bar.c   (working copy)
@@ -1,7 +1,12 @@
+#include <sys/types.h>
+#include <sys/stat.h>
+#include <unistd.h>
+
+#include <stdio.h>

 int main(void) {
-  printf("Sixty-four slices of American Cheese...\n");
+  printf("Sixty-five slices of American Cheese...\n");
 return 0;
 }

Index: README
===================================================================
--- README  (revision 3)
+++ README  (working copy)
@@ -193,3 +193,4 @@
+Note to self:  pick up laundry.

Index: stuff/fish.c
===================================================================
--- stuff/fish.c    (revision 1)
+++ stuff/fish.c    (working copy)
-Welcome to the file known as 'fish'.
-Information on fish will be here soon.

Index: stuff/things/bloo.h
===================================================================
--- stuff/things/bloo.h (revision 8)
+++ stuff/things/bloo.h (working copy)
+Here is a new file to describe
+things about bloo.
\end{verbatim}

The \texttt{svn\ diff} command produces this output by comparing your
working files against its pristine text-base. Files scheduled for
addition are displayed as files in which every line was added; files
scheduled for deletion are displayed as if every line was removed from
those files. The output from \texttt{svn\ diff} is somewhat compatible
with the \texttt{patch} program---more so with the \texttt{svn\ patch}
subcommand introduced in Subversion 1.7. Patch processing commands such
as these read and apply patch files (or ``patches''), which are files
that describe differences made to one or more files. Because of this,
you can share the changes you\textquotesingle ve made in your working
copy with someone else without first committing those changes by
creating a patch file from the redirected output of \texttt{svn\ diff}:

\begin{verbatim}
$ svn diff > patchfile
$
\end{verbatim}

Subversion uses its internal diff engine, which produces unified diff
format, by default. If you want diff output in a different format,
specify an external diff program using \texttt{-\/-diff-cmd} and pass
any additional flags that it needs via the \texttt{-\/-extensions}
(\texttt{-x}) option. For example, you might want Subversion to defer
its difference calculation and display to the GNU \texttt{diff} program,
asking that program to print local modifications made to the file
\texttt{foo.c} in context diff format (another flavor of difference
format) while ignoring changes made only to the case of the letters used
in the file\textquotesingle s contents:

\begin{verbatim}
$ svn diff --diff-cmd /usr/bin/diff -x "-i" foo.c
…
$
\end{verbatim}

\subsection{Fix Your Mistakes}\label{svn.tour.cycle.revert}

Suppose while viewing the output of \texttt{svn\ diff} you determine
that all the changes you made to a particular file are mistakes. Maybe
you shouldn\textquotesingle t have changed the file at all, or perhaps
it would be easier to make different changes starting from scratch. You
could edit the file again and unmake all those changes. You could try to
find a copy of how the file looked before you changed it, and then copy
its contents atop your modified version. You could attempt to apply
those changes to the file again in reverse using
\texttt{svn\ patch\ -\/-reverse-diff} or using your operating
system\textquotesingle s \texttt{patch\ -R}. And there are probably
other approaches you could take.

Fortunately in Subversion, undoing your work and starting over from
scratch doesn\textquotesingle t require such acrobatics. Just use the
\texttt{svn\ revert} command:

\begin{verbatim}
$ svn status README
M       README
$ svn revert README
Reverted 'README'
$ svn status README
$
\end{verbatim}

In this example, Subversion has reverted the file to its premodified
state by overwriting it with the pristine version of the file cached in
the text-base area. But note that \texttt{svn\ revert} can undo
\emph{any} scheduled operation---for example, you might decide that you
don\textquotesingle t want to add a new file after all:

\begin{verbatim}
$ svn status new-file.txt
?       new-file.txt
$ svn add new-file.txt
A         new-file.txt
$ svn revert new-file.txt
Reverted 'new-file.txt'
$ svn status new-file.txt
?       new-file.txt
$
\end{verbatim}

Or perhaps you mistakenly removed a file from version control:

\begin{verbatim}
$ svn status README
$ svn delete README
D         README
$ svn revert README
Reverted 'README'
$ svn status README
$
\end{verbatim}

The \texttt{svn\ revert} command offers salvation for imperfect people.
It can save you huge amounts of time and energy that would otherwise be
spent manually unmaking changes or, worse, disposing of your working
copy and checking out a fresh one just to have a clean slate to work
with again.

\subsection{Resolve Any Conflicts}\label{svn.tour.cycle.resolve}

We\textquotesingle ve already seen how \texttt{svn\ status\ -u} can
predict conflicts, but dealing with those conflicts is still something
that remains to be done. Conflicts can occur any time you attempt to
merge or integrate (in a very general sense) changes from the repository
into your working copy. By now you know that \texttt{svn\ update}
creates exactly that sort of scenario---that command\textquotesingle s
very purpose is to bring your working copy up to date with the
repository by merging all the changes made since your last update into
your working copy. So how does Subversion report these conflicts to you,
and how do you deal with them?

Suppose you run \texttt{svn\ update} and you see this sort of
interesting output:

\begin{verbatim}
$ svn update
Updating '.':
U    INSTALL
G    README
Conflict discovered in 'bar.c'.
Select: (p) postpone, (df) show diff, (e) edit file, (m) merge,
        (mc) my side of conflict, (tc) their side of conflict,
        (s) show all options:
        
\end{verbatim}

The \texttt{U} (which stands for ``Updated'') and \texttt{G} (for
``merGed'') codes are no cause for concern; those files cleanly absorbed
changes from the repository. A file marked with \texttt{U} contains no
local changes but was updated with changes from the repository. One
marked with \texttt{G} had local changes to begin with, but the changes
coming from the repository didn\textquotesingle t conflict with those
local changes.

It\textquotesingle s the next few lines which are interesting. First,
Subversion reports to you that in its attempt to merge outstanding
server changes into the file \texttt{bar.c}, it has detected that some
of those changes clash with local modifications you\textquotesingle ve
made to that file in your working copy but have not yet committed.
Perhaps someone has changed the same line of text you also changed.
Whatever the reason, Subversion instantly flags this file as being in a
state of conflict. It then asks you what you want to do about the
problem, allowing you to interactively choose an action to take toward
resolving the conflict. The most commonly used options are displayed,
but you can see all of the options by typing \textless s\textgreater:

\begin{verbatim}
…
Select: (p) postpone, (df) show diff, (e) edit file, (m) merge,
        (mc) my side of conflict, (tc) their side of conflict,
        (s) show all options: s

  (e)  - change merged file in an editor  [edit]
  (df) - show all changes made to merged file
  (r)  - accept merged version of file

  (dc) - show all conflicts (ignoring merged version)
  (mc) - accept my version for all conflicts (same)  [mine-conflict]
  (tc) - accept their version for all conflicts (same)  [theirs-conflict]

  (mf) - accept my version of entire file (even non-conflicts)  [mine-full]
  (tf) - accept their version of entire file (same)  [theirs-full]

  (m)  - use internal merge tool to resolve conflict
  (l)  - launch external tool to resolve conflict  [launch]
  (p)  - mark the conflict to be resolved later  [postpone]
  (q)  - postpone all remaining conflicts
  (s)  - show this list (also 'h', '?')
Words in square brackets are the corresponding --accept option arguments.

Select: (p) postpone, (df) show diff, (e) edit file, (m) merge,
        (mc) my side of conflict, (tc) their side of conflict,
        (s) show all options:
\end{verbatim}

Let\textquotesingle s briefly review each of these options before we go
into detail on what each option means.

\begin{description}
\item[\texttt{(e)\ edit\ {[}edit{]}}]
Open the file in conflict with your favorite editor, as set in the
environment variable \texttt{EDITOR}.
\item[\texttt{(df)\ diff-full}]
Display the differences between the base revision and the conflicted
file itself in unified diff format.
\item[\texttt{(r)\ resolved}]
After editing a file, tell \texttt{svn} that you\textquotesingle ve
resolved the conflicts in the file and that it should accept the current
contents---basically that you\textquotesingle ve ``resolved'' the
conflict.
\item[\texttt{(dc)\ display-conflict}]
Display all conflicting regions of the file, ignoring changes which were
successfully merged.
\item[\texttt{(mc)\ mine-conflict\ {[}mine-conflict{]}}]
Discard any newly received changes from the server which conflict with
your local changes to the file under review. However, accept and merge
all non-conflicting changes received from the server for that file.
\item[\texttt{(tc)\ theirs-conflict\ {[}theirs-conflict{]}}]
Discard any local changes which conflict with incoming changes from the
server for the file under review. However, preserve all non-conflicting
local changes to that file.
\item[\texttt{(mf)\ mine-full\ {[}mine-full{]}}]
Discard all newly received changes from the server for the file under
review, but preserve all your local changes for that file.
\item[\texttt{(tf)\ theirs-full\ {[}theirs-full{]}}]
Discard all your local changes to the file under review and use only the
newly received changes from the server for that file.
\item[\texttt{(m)\ merge}]
Launch an internal file merge tool to perform the conflict resolution.
The option is available starting with Subversion 1.8.
\item[\texttt{(l)\ launch}]
Launch an external program to perform the conflict resolution. This
requires a bit of preparation beforehand.
\item[\texttt{(p)\ postpone\ {[}postpone{]}}]
Leave the file in a conflicted state for you to resolve after your
update is complete.
\item[\texttt{(s)\ show\ all}]
Show the list of all possible commands you can use in interactive
conflict resolution.
\end{description}

We\textquotesingle ll cover these commands in more detail now, grouping
them together by related functionality.

\subsubsection{Viewing conflict differences
interactively}\label{svn.tour.cycle.resolve.diff}

Before deciding how to attack a conflict interactively, odds are that
you\textquotesingle d like to see exactly what is in conflict. Two of
the commands available at the interactive conflict resolution prompt can
assist you here. The first is the ``diff-full'' command (\texttt{df}),
which displays all the local modifications to the file in question plus
any conflict regions:

\begin{verbatim}
…
Select: (p) postpone, (df) show diff, (e) edit file, (m) merge,
        (mc) my side of conflict, (tc) their side of conflict,
        (s) show all options: df
--- .svn/text-base/sandwich.txt.svn-base      Tue Dec 11 21:33:57 2007
+++ .svn/tmp/tempfile.32.tmp     Tue Dec 11 21:34:33 2007
@@ -1 +1,5 @@
-Just buy a sandwich.
+<<<<<<< .mine
+Go pick up a cheesesteak.
+=======
+Bring me a taco!
+>>>>>>> .r32
…
\end{verbatim}

The first line of the diff content shows the previous contents of the
working copy (the \texttt{BASE} revision), the next content line is your
change, and the last content line is the change that was just received
from the server (\emph{usually} the \texttt{HEAD} revision).

The second command is similar to the first, but the ``display-conflict''
(\texttt{dc}) command shows only the conflict regions, not all the
changes made to the file. Additionally, this command uses a slightly
different display format for the conflict regions which allows you to
more easily compare the file\textquotesingle s contents in those regions
as they would appear in each of three states: original and unedited;
with your local changes applied and the server\textquotesingle s
conflicting changes ignored; and with only the server\textquotesingle s
incoming changes applied and your local, conflicting changes reverted.

After reviewing the information provided by these commands,
you\textquotesingle re ready to move on to the next action.

\subsubsection{Resolving conflict differences
interactively}\label{svn.tour.cycle.resolve.resolve}

The main way to resolve conflicts interactively is to use an internal
file merge tool. The tool asks you what to do with each conflicting
change and allows you to selectively merge and edit changes. However,
there are several other different ways to resolve conflicts
interactively---two of them allow you to selectively merge and edit
changes using external editors, the rest of which allow you to simply
pick a version of the file and move along. Internal merge tool combines
all of the available ways to resolve conflicts.

You\textquotesingle ve already reviewed the conflicting changes, so
it\textquotesingle s now time to resolve the conflicts. The first
command that should help you is the ``merge'' command (\texttt{m}) which
is available starting with Subversion 1.8. The command displays the
conflicting areas and allows you to choose from a number of options to
resolve the conflicts area-by-area:

\begin{verbatim}
Select: (p) postpone, (df) show diff, (e) edit file, (m) merge,
        (mc) my side of conflict, (tc) their side of conflict,
        (s) show all options: m
Merging 'Makefile'.
Conflicting section found during merge:
(1) their version (at line 24)                  |(2) your version (at line 24)
------------------------------------------------+------------------------------------------------
top_builddir = /bar                             |top_builddir = /foo
------------------------------------------------+------------------------------------------------
Select: (1) use their version, (2) use your version,
        (12) their version first, then yours,
        (21) your version first, then theirs,
        (e1) edit their version and use the result,
        (e2) edit your version and use the result,
        (eb) edit both versions and use the result,
        (p) postpone this conflicting section leaving conflict markers,
        (a) abort file merge and return to main menu:
\end{verbatim}

As you can see, when you use the internal file merge tool, you can cycle
through individual conflicting areas in the file and select various
resolution options or postpone conflict resolution for selected
conflicts.

However, if you wish to use an external editor to choose some
combination of your local changes, you can use the ``edit'' command
(\texttt{e}) to manually edit the file with conflict markers in a text
editor (configured per the instructions in
\hyperref[svn.advanced.externaleditors]{Using External Editors}). After
you\textquotesingle ve edited the file, if you\textquotesingle re
satisfied with the changes you\textquotesingle ve made, you can tell
Subversion that the edited file is no longer in conflict by using the
``resolved'' command (\texttt{r}).

Regardless of what your local Unix snob will likely tell you, editing
the file by hand in your favorite text editor is a somewhat low-tech way
of remedying conflicts (see
\hyperref[svn.tour.cycle.resolve.byhand]{Manual conflict resolution} for
a walkthrough). For this reason, Subversion provides the ``launch''
resolution command (\texttt{l}) to fire up a fancy graphical merge tool
instead (see \hyperref[svn.advanced.externaldifftools.merge]{External
merge}).

There is also a pair of compromise options available. The
``mine-conflict'' (\texttt{mc}) and ``theirs-conflict'' (\texttt{tc})
commands instruct Subversion to select your local changes or the
server\textquotesingle s incoming changes, respectively, as the
``winner'' for all conflicts in the file. But, unlike the ``mine-full''
and ``theirs-full'' commands, these commands preserve both your local
changes and changes received from the server in regions of the file
where no conflict was detected.

Finally, if you decide that you don\textquotesingle t need to merge any
changes, but just want to accept one version of the file or the other,
you can either choose your changes (a.k.a. ``mine'') by using the
``mine-full'' command (\texttt{mf}) or choose theirs by using the
``theirs-full'' command (\texttt{tf}).

\subsubsection{Postponing conflict
resolution}\label{svn.tour.cycle.resolve.pending}

This may sound like an appropriate section for avoiding marital
disagreements, but it\textquotesingle s actually still about Subversion,
so read on. If you\textquotesingle re doing an update and encounter a
conflict that you\textquotesingle re not prepared to review or resolve,
you can type \texttt{p} to postpone resolving a conflict on a
file-by-file basis when you run \texttt{svn\ update}. If you know in
advance that you don\textquotesingle t want to resolve any conflicts
interactively, you can pass the \texttt{-\/-non-interactive} option to
\texttt{svn\ update}, and any file in conflict will be marked with a
\texttt{C} automatically.

Beginning with Subversion 1.8, an internal file merge tool allows you to
postpone conflict resolution for certain conflicts, but resolve other
conflicts. Therefore, you can postpone conflict resolution area-by-area,
not just on a file-to-file basis.

The \texttt{C} (for ``Conflicted'') means that the changes from the
server overlapped with your own, and now you have to manually choose
between them after the update has completed. When you postpone a
conflict resolution, \texttt{svn} typically does three things to assist
you in noticing and resolving that conflict:

\begin{itemize}
\item
  Subversion prints a \texttt{C} during the update and remembers that
  the file is in a state of conflict.
\item
  If Subversion considers the file to be mergeable, it places conflict
  markers---special strings of text that delimit the ``sides'' of the
  conflict---into the file to visibly demonstrate the overlapping areas.
  (Subversion uses the \texttt{svn:mime-type} property to decide whether
  a file is capable of contextual, line-based merging. See
  \hyperref[svn.advanced.props.special.mime-type]{File Content Type} to
  learn more.)
\item
  For every conflicted file, Subversion places three extra unversioned
  files in your working copy:

  \begin{description}
  \item[\texttt{filename.mine}]
  This is the file as it existed in your working copy before you began
  the update process. It contains any local modifications you had made
  to the file up to that point. (If Subversion considers the file to be
  unmergeable, the \texttt{.mine} file isn\textquotesingle t created,
  since it would be identical to the working file.)
  \item[\texttt{filename.rOLDREV}]
  This is the file as it existed in the \texttt{BASE} revision---that
  is, the unmodified revision of the file in your working copy
  \emph{before} you began the update process---where
  \textless OLDREV\textgreater{} is that base revision number.
  \item[\texttt{filename.rNEWREV}]
  This is the file that your Subversion client just received from the
  server via the update of your working copy, where
  \textless NEWREV\textgreater{} corresponds to the revision number to
  which you were updating (\texttt{HEAD}, unless otherwise requested).
  \end{description}
\end{itemize}

For example, Sally makes changes to the file \texttt{sandwich.txt}, but
does not yet commit those changes. Meanwhile, Harry commits changes to
that same file. Sally updates her working copy before committing and she
gets a conflict, which she postpones:

\begin{verbatim}
$ svn update
Updating '.':
Conflict discovered in 'sandwich.txt'.
Select: (p) postpone, (df) show diff, (e) edit file, (m) merge,
        (mc) my side of conflict, (tc) their side of conflict,
        (s) show all options: p
C    sandwich.txt
Updated to revision 2.
Summary of conflicts:
  Text conflicts: 1
$ ls -1
sandwich.txt
sandwich.txt.mine
sandwich.txt.r1
sandwich.txt.r2
\end{verbatim}

At this point, Subversion will \emph{not} allow Sally to commit the file
\texttt{sandwich.txt} until the three temporary files are removed:

\begin{verbatim}
$ svn commit -m "Add a few more things"
svn: E155015: Commit failed (details follow):
svn: E155015: Aborting commit: '/home/sally/svn-work/sandwich.txt' remains in conflict
\end{verbatim}

If you\textquotesingle ve postponed a conflict, you need to resolve the
conflict before Subversion will allow you to commit your changes.
You\textquotesingle ll do this with the \texttt{svn\ resolve} command.
This command accepts the \texttt{-\/-accept} option, which allows you
specify your desired approach for resolving the conflict. Prior to
Subversion 1.8, the \texttt{svn\ resolve} command \emph{required} the
use of this option. Subversion now allows you to run the
\texttt{svn\ resolve} command without that option. When you do so,
Subversion cranks up its interactive conflict resolution mechanism,
which you can read about (if you haven\textquotesingle t done so
already) in the previous section,
\hyperref[svn.tour.cycle.resolve.resolve]{Resolving conflict differences
interactively}. We\textquotesingle ll take the opportunity in this
section, though, to discuss the use of the \texttt{-\/-accept} option
for conflict resolution.

The \texttt{-\/-accept} option to the \texttt{svn\ resolve} command
instructs Subversion to use one of its pre-packaged approaches to
conflict resolution. If you want Subversion to resolve the conflict
using the version of the file that you last checked out before making
your edits, use \texttt{-\/-accept=base}. If you\textquotesingle d
prefer instead to keep the version that contains only your edits, use
\texttt{-\/-accept=mine-full}. You can also select the version that your
most recent update pulled from the server (discarding your edits
entirely)---that\textquotesingle s done using
\texttt{-\/-accept=theirs-full}. There are other ``canned'' resolution
types, too. See \hyperref[svn.ref.svn.sw.accept]{\texttt{-\/-accept}
\textless ACTION\textgreater{}} in \hyperref[svn.ref.svn]{svn
Reference---Subversion Command-Line Client} for details.

You aren\textquotesingle t limited strictly to all-or-nothing options.
If you want to pick and choose from your changes and the changes that
your update fetched from the server, you can manually repair the working
file, fixing up the conflicted text ``by hand'' (by examining and
editing the conflict markers within the file), then tell Subversion to
resolve the conflict by keeping the working file in its current state by
running \texttt{svn\ resolve} with the \texttt{-\/-accept=working}
option.

\texttt{svn\ resolve} removes the three temporary files and accepts the
version of the file that you specified. After the command completes
successfully---and assuming you didn\textquotesingle t interactively
choose to postpone resolution, of course---Subversion no longer
considers the file to be in a state of conflict:

\begin{verbatim}
$ svn resolve --accept working sandwich.txt
Resolved conflicted state of 'sandwich.txt'
\end{verbatim}

\subsubsection{Manual conflict
resolution}\label{svn.tour.cycle.resolve.byhand}

Manually resolving conflicts can be quite intimidating the first time
you attempt it, but with a little practice, it can become as easy as
falling off a bike.

Here\textquotesingle s an example. Due to a miscommunication, you and
Sally, your collaborator, both edit the file \texttt{sandwich.txt} at
the same time. Sally commits her changes, and when you go to update your
working copy, you get a conflict and you\textquotesingle re going to
have to edit \texttt{sandwich.txt} to resolve the conflict. First,
let\textquotesingle s take a look at the file:

\begin{verbatim}
$ cat sandwich.txt
Top piece of bread
Mayonnaise
Lettuce
Tomato
Provolone
<<<<<<< .mine
Salami
Mortadella
Prosciutto
=======
Sauerkraut
Grilled Chicken
>>>>>>> .r2
Creole Mustard
Bottom piece of bread
\end{verbatim}

The strings of less-than signs, equals signs, and greater-than signs are
conflict markers and are not part of the actual data in conflict. You
generally want to ensure that those are removed from the file before
your next commit. The text between the first two sets of markers is
composed of the changes you made in the conflicting area:

\begin{verbatim}
<<<<<<< .mine
Salami
Mortadella
Prosciutto
=======
\end{verbatim}

The text between the second and third sets of conflict markers is the
text from Sally\textquotesingle s commit:

\begin{verbatim}
=======
Sauerkraut
Grilled Chicken
>>>>>>> .r2
\end{verbatim}

Usually you won\textquotesingle t want to just delete the conflict
markers and Sally\textquotesingle s changes---she\textquotesingle s
going to be awfully surprised when the sandwich arrives and
it\textquotesingle s not what she wanted. This is where you pick up the
phone or walk across the office and explain to Sally that you
can\textquotesingle t get sauerkraut from an Italian deli.\footnote{And
  if you ask them for it, they may very well ride you out of town on a
  rail.} Once you\textquotesingle ve agreed on the changes you will
commit, edit your file and remove the conflict markers:

\begin{verbatim}
Top piece of bread
Mayonnaise
Lettuce
Tomato
Provolone
Salami
Mortadella
Prosciutto
Creole Mustard
Bottom piece of bread
\end{verbatim}

Now use \texttt{svn\ resolve}, and you\textquotesingle re ready to
commit your changes:

\begin{verbatim}
$ svn resolve --accept working sandwich.txt
Resolved conflicted state of 'sandwich.txt'
$ svn commit -m "Go ahead and use my sandwich, discarding Sally's edits."
\end{verbatim}

Naturally, you want to be careful that when using \texttt{svn\ resolve}
you don\textquotesingle t tell Subversion that you\textquotesingle ve
resolved a conflict when you truly haven\textquotesingle t. Once the
temporary files are removed, Subversion will let you commit the file
even if it still contains conflict markers.

If you ever get confused while editing the conflicted file, you can
always consult the three files that Subversion creates for you in your
working copy---including your file as it was before you updated. You can
even use a third-party interactive merging tool to examine those three
files.

\subsubsection{Discarding your changes in favor of a newly fetched
revision}\label{svn.tour.cycle.resolve.theirsfull}

If you get a conflict and decide that you want to throw out your
changes, you can run
\texttt{svn\ resolve\ -\/-accept\ theirs-full\ CONFLICTED-PATH} and
Subversion will discard your edits and remove the temporary files:

\begin{verbatim}
$ svn update
Updating '.':
Conflict discovered in 'sandwich.txt'.
Select: (p) postpone, (df) show diff, (e) edit file, (m) merge,
        (mc) my side of conflict, (tc) their side of conflict,
        (s) show all options: p
C    sandwich.txt
Updated to revision 2.
Summary of conflicts:
  Text conflicts: 1
$ ls sandwich.*
sandwich.txt  sandwich.txt.mine  sandwich.txt.r2  sandwich.txt.r1
$ svn resolve --accept theirs-full sandwich.txt
Resolved conflicted state of 'sandwich.txt'
$
\end{verbatim}

\subsubsection{Punting: using svn
revert}\label{svn.tour.cycle.resolve.revert}

If you decide that you want to throw out your changes and start your
edits again (whether this occurs after a conflict or anytime), just
revert your changes:

\begin{verbatim}
$ svn revert sandwich.txt
Reverted 'sandwich.txt'
$ ls sandwich.*
sandwich.txt
$
\end{verbatim}

Note that when you revert a conflicted file, you don\textquotesingle t
have to use \texttt{svn\ resolve}.

\subsection{Commit Your Changes}\label{svn.tour.cycle.commit}

Finally! Your edits are finished, you\textquotesingle ve merged all
changes from the server, and you\textquotesingle re ready to commit your
changes to the repository.

The \texttt{svn\ commit} command sends all of your changes to the
repository. When you commit a change, you need to supply a log message
describing your change. Your log message will be attached to the new
revision you create. If your log message is brief, you may wish to
supply it on the command line using the \texttt{-\/-message}
(\texttt{-m}) option:

\begin{verbatim}
$ svn commit -m "Corrected number of cheese slices."
Sending        sandwich.txt
Transmitting file data .
Committed revision 3.
\end{verbatim}

However, if you\textquotesingle ve been composing your log message in
some other text file as you work, you may want to tell Subversion to get
the message from that file by passing its filename as the value of the
\texttt{-\/-file} (\texttt{-F}) option:

\begin{verbatim}
$ svn commit -F logmsg
Sending        sandwich.txt
Transmitting file data .
Committed revision 4.
\end{verbatim}

If you fail to specify either the \texttt{-\/-message} (\texttt{-m}) or
\texttt{-\/-file} (\texttt{-F}) option, Subversion will automatically
launch your favorite editor (see the information on \texttt{editor-cmd}
in \hyperref[svn.advanced.confarea.opts.config]{General configuration})
for composing a log message.

If you\textquotesingle re in your editor writing a commit message and
decide that you want to cancel your commit, you can just quit your
editor without saving changes. If you\textquotesingle ve already saved
your commit message, simply delete all the text, save again, and then
abort:

\begin{verbatim}
$ svn commit
Waiting for Emacs...Done

Log message unchanged or not specified
(a)bort, (c)ontinue, (e)dit
a
$
\end{verbatim}

The repository doesn\textquotesingle t know or care whether your changes
make any sense as a whole; it checks only to make sure nobody else has
changed any of the same files that you did when you
weren\textquotesingle t looking. If somebody \emph{has} done that, the
entire commit will fail with a message informing you that one or more of
your files are out of date:

\begin{verbatim}
$ svn commit -m "Add another rule"
Sending        rules.txt
Transmitting file data .
svn: E155011: Commit failed (details follow):
svn: E155011: File '/home/sally/svn-work/sandwich.txt' is out of date
…
\end{verbatim}

(The exact wording of this error message depends on the network protocol
and server you\textquotesingle re using, but the idea is the same in all
cases.)

At this point, you need to run \texttt{svn\ update}, deal with any
merges or conflicts that result, and then attempt your commit again.

That covers the basic work cycle for using Subversion. Subversion offers
many other features that you can use to manage your repository and
working copy, but most of your day-to-day use of Subversion will involve
only the commands that we\textquotesingle ve discussed so far in this
chapter. We will, however, cover a few more commands that
you\textquotesingle ll use fairly often.

\section{Examining History}\label{svn.tour.history}

Your Subversion repository is like a time machine. It keeps a record of
every change ever committed and allows you to explore this history by
examining previous versions of files and directories as well as the
metadata that accompanies them. With a single Subversion command, you
can check out the repository (or restore an existing working copy)
exactly as it was at any date or revision number in the past. However,
sometimes you just want to \emph{peer into} the past instead of
\emph{going into} it.

Several commands can provide you with historical data from the
repository:

\begin{description}
\item[\texttt{svn\ diff}]
Shows line-level details of a particular change
\item[\texttt{svn\ log}]
Shows you broad information: log messages with date and author
information attached to revisions and which paths changed in each
revision
\item[\texttt{svn\ cat}]
Retrieves a file as it existed in a particular revision number and
displays it on your screen
\item[\texttt{svn\ annotate}]
Retrieves a human-readable file as it existed in a particular revision
number, displaying its contents in a tabular form with last-changed
information added to each line of the file
\item[\texttt{svn\ list}]
Displays the files in a directory for any given revision
\end{description}

\subsection{Examining the Details of Historical
Changes}\label{svn.tour.history.diff}

We\textquotesingle ve already seen \texttt{svn\ diff} before---it
displays file differences in unified diff format; we used it to show the
local modifications made to our working copy before committing to the
repository.

In fact, it turns out that there are \emph{three} distinct uses of
\texttt{svn\ diff}:

\begin{itemize}
\item
  Examining local changes
\item
  Comparing your working copy to the repository
\item
  Comparing repository revisions
\end{itemize}

\subsubsection{Examining local
changes}\label{svn.tour.history.diff.local}

As we\textquotesingle ve seen, invoking \texttt{svn\ diff} with no
options will compare your working files to the cached ``pristine''
copies in the \texttt{.svn} area:

\begin{verbatim}
$ svn diff
Index: rules.txt
===================================================================
--- rules.txt   (revision 3)
+++ rules.txt   (working copy)
@@ -1,4 +1,5 @@
 Be kind to others
 Freedom = Responsibility
 Everything in moderation
-Chew with your mouth open
+Chew with your mouth closed
+Listen when others are speaking
$
\end{verbatim}

\subsubsection{Comparing working copy to
repository}\label{svn.tour.history.diff.wcrepos}

If a single \texttt{-\/-revision} (\texttt{-r}) number is passed, your
working copy is compared to the specified revision in the repository:

\begin{verbatim}
$ svn diff -r 3 rules.txt
Index: rules.txt
===================================================================
--- rules.txt   (revision 3)
+++ rules.txt   (working copy)
@@ -1,4 +1,5 @@
 Be kind to others
 Freedom = Responsibility
 Everything in moderation
-Chew with your mouth open
+Chew with your mouth closed
+Listen when others are speaking
$
\end{verbatim}

\subsubsection{Comparing repository
revisions}\label{svn.tour.history.diff.reposrepos}

If two revision numbers, separated by a colon, are passed via
\texttt{-\/-revision} (\texttt{-r}), the two revisions are directly
compared:

\begin{verbatim}
$ svn diff -r 2:3 rules.txt
Index: rules.txt
===================================================================
--- rules.txt   (revision 2)
+++ rules.txt   (revision 3)
@@ -1,4 +1,4 @@
 Be kind to others
-Freedom = Chocolate Ice Cream
+Freedom = Responsibility
 Everything in moderation
 Chew with your mouth open
$
\end{verbatim}

A more convenient way of comparing one revision to the previous revision
is to use the \texttt{-\/-change} (\texttt{-c}) option:

\begin{verbatim}
$ svn diff -c 3 rules.txt
Index: rules.txt
===================================================================
--- rules.txt   (revision 2)
+++ rules.txt   (revision 3)
@@ -1,4 +1,4 @@
 Be kind to others
-Freedom = Chocolate Ice Cream
+Freedom = Responsibility
 Everything in moderation
 Chew with your mouth open
$
\end{verbatim}

Lastly, you can compare repository revisions even when you
don\textquotesingle t have a working copy on your local machine, just by
including the appropriate URL on the command line:

\begin{verbatim}
$ svn diff -c 5 http://svn.example.com/repos/example/trunk/text/rules.txt
…
$
\end{verbatim}

\subsection{Generating a List of Historical
Changes}\label{svn.tour.history.log}

To find information about the history of a file or directory, use the
\texttt{svn\ log} command. \texttt{svn\ log} will provide you with a
record of who made changes to a file or directory, at what revision it
changed, the time and date of that revision, and---if it was
provided---the log message that accompanied the commit:

\begin{verbatim}
$ svn log
------------------------------------------------------------------------
r3 | sally | 2008-05-15 23:09:28 -0500 (Thu, 15 May 2008) | 1 line

Added include lines and corrected # of cheese slices.
------------------------------------------------------------------------
r2 | harry | 2008-05-14 18:43:15 -0500 (Wed, 14 May 2008) | 1 line

Added main() methods.
------------------------------------------------------------------------
r1 | sally | 2008-05-10 19:50:31 -0500 (Sat, 10 May 2008) | 1 line

Initial import
------------------------------------------------------------------------
\end{verbatim}

Note that the log messages are printed in \emph{reverse chronological
order} by default. If you wish to see a different range of revisions in
a particular order or just a single revision, pass the
\texttt{-\/-revision} (\texttt{-r}) option:

\begin{longtable}[]{@{}>{\raggedright\arraybackslash}p{\dimexpr(\linewidth-2\tabcolsep*2)/2\relax}>{\raggedright\arraybackslash}p{\dimexpr(\linewidth-2\tabcolsep*2)/2\relax}@{}}
\caption{Common log
requests}\label{svn.tour.history.log.tbl-1}\tabularnewline
\toprule\noalign{}
Command & Description \\
\midrule\noalign{}
\endfirsthead
\toprule\noalign{}
Command & Description \\
\midrule\noalign{}
\endhead
\bottomrule\noalign{}
\endlastfoot
\texttt{svn\ log\ -r\ 5:19} & Display logs for revisions 5 through 19 in
chronological order \\
\texttt{svn\ log\ -r\ 19:5} & Display logs for revisions 5 through 19 in
reverse chronological order \\
\texttt{svn\ log\ -r\ 8} & Display logs for revision 8 only \\
\end{longtable}

You can also examine the log history of a single file or directory. For
example:

\begin{verbatim}
$ svn log foo.c
…
$ svn log http://foo.com/svn/trunk/code/foo.c
…
\end{verbatim}

These will display log messages \emph{only} for those revisions in which
the named file (or directory) changed.

Why Does svn log Not Show Me What I Just Committed?

If you make a commit and immediately type \texttt{svn\ log} with no
arguments, you may notice that your most recent commit
doesn\textquotesingle t show up in the list of log messages. This is due
to a combination of the behavior of \texttt{svn\ commit} and the default
behavior of \texttt{svn\ log}. First, when you commit changes to the
repository, \texttt{svn} bumps only the revision number of files (and
directories) that it commits, so usually the parent directory remains at
the older revision (See
\hyperref[svn.basic.in-action.mixedrevs.update-commit]{Updates and
commits are separate} for an explanation of why). \texttt{svn\ log} then
defaults to fetching the history of the directory at its current
revision, and thus you don\textquotesingle t see the newly committed
changes. The solution here is to either update your working copy or
explicitly provide a revision number to \texttt{svn\ log} by using the
\texttt{-\/-revision} (\texttt{-r}) option.

If you want even more information about a file or directory,
\texttt{svn\ log} also takes a \texttt{-\/-verbose} (\texttt{-v})
option. Because Subversion allows you to move and copy files and
directories, it is important to be able to track path changes in the
filesystem. So, in verbose mode, \texttt{svn\ log} will include a list
of changed paths in a revision in its output:

\begin{verbatim}
$ svn log -r 8 -v
------------------------------------------------------------------------
r8 | sally | 2008-05-21 13:19:25 -0500 (Wed, 21 May 2008) | 1 line
Changed paths:
   M /trunk/code/foo.c
   M /trunk/code/bar.h
   A /trunk/code/doc/README

Frozzled the sub-space winch.

------------------------------------------------------------------------
\end{verbatim}

\texttt{svn\ log} also takes a \texttt{-\/-quiet} (\texttt{-q}) option,
which suppresses the body of the log message. When combined with
\texttt{-\/-verbose} (\texttt{-v}), it gives just the names of the
changed files.

Why Does svn log Give Me an Empty Response?

After working with Subversion for a bit, most users will come across
something like this:

\begin{verbatim}
$ svn log -r 2
------------------------------------------------------------------------
$
\end{verbatim}

At first glance, this seems like an error. But recall that while
revisions are repository-wide, \texttt{svn\ log} operates on a path in
the repository. If you supply no path, Subversion uses the current
working directory as the default target. As a result, if
you\textquotesingle re operating in a subdirectory of your working copy
and attempt to see the log of a revision in which neither that directory
nor any of its children was changed, Subversion will show you an empty
log. If you want to see what changed in that revision, try pointing
\texttt{svn\ log} directly at the topmost URL of your repository, as in
\texttt{svn\ log\ -r\ 2\ \^{}/}.

As of Subversion 1.7, users of the Subversion command-line can also take
advantage of a special output mode for \texttt{svn\ log} which
integrates a difference report such as is generated by the
\texttt{svn\ diff} command we introduced earlier. When you invoke
\texttt{svn\ log} with the \texttt{-\/-diff} option, Subversion will
append to each revision log chunk in the log report a
\texttt{diff}-style difference report. This is a very convenient way to
see both the high-level, semantic changes and the line-based
modifications of a revision all at the same time!

Beginning with Subversion 1.8, \texttt{svn\ log} accepts
\texttt{-\/-search} and \texttt{-\/-search-and} options. The options
allow you to filter the output of \texttt{svn\ log} based on the search
pattern you supply. When using these options, a log message is shown
only if a revision\textquotesingle s author, date, log message text, or
list of changed paths, matches the search pattern.

\subsection{Browsing the Repository}\label{svn.tour.history.browsing}

Using \texttt{svn\ cat} and \texttt{svn\ list}, you can view various
revisions of files and directories without changing the working revision
of your working copy. In fact, you don\textquotesingle t even need a
working copy to use either one.

\subsubsection{Displaying file
contents}\label{svn.tour.history.browsing.cat}

If you want to examine an earlier version of a file and not necessarily
the differences between two files, you can use \texttt{svn\ cat}:

\begin{verbatim}
$ svn cat -r 2 rules.txt
Be kind to others
Freedom = Chocolate Ice Cream
Everything in moderation
Chew with your mouth open
$
\end{verbatim}

You can also redirect the output directly into a file:

\begin{verbatim}
$ svn cat -r 2 rules.txt > rules.txt.v2
$
\end{verbatim}

\subsubsection{Displaying line-by-line change
attribution}\label{svn.tour.history.browsing.annotate}

Very similar to the \texttt{svn\ cat} command we discussed in the
previous section is the \texttt{svn\ annotate} command. This command
also displays the contents of a versioned file, but it does so using a
tabular format. Each line of output shows not only a line of the
file\textquotesingle s content but also the username, the revision
number and (optionally) the datestamp of the revision in which that line
was last modified.

When used with a working copy file target, \texttt{svn\ annotate} will
by default show line-by-line attribution of the file as it currently
appears in the working copy.

\begin{verbatim}
$ svn annotate rules.txt
     1      harry Be kind to others
     3      sally Freedom = Responsibility
     1      harry Everything in moderation
     -          - Chew with your mouth closed
     -          - Listen when others are speaking
\end{verbatim}

Notice that for some lines, there is no attribution provided. In this
case, that\textquotesingle s because those lines have been modified in
the working copy\textquotesingle s version of the file. In this way,
\texttt{svn\ annotate} becomes another way for you to see which lines in
the file you have changed. You can use the \texttt{BASE} revision
keyword (see \hyperref[svn.tour.revs.keywords]{Revision Keywords}) to
instead see the unmodified form of the file as it resides in your
working copy.

\begin{verbatim}
$ svn annotate rules.txt@BASE
     1      harry Be kind to others
     3      sally Freedom = Responsibility
     1      harry Everything in moderation
     1      harry Chew with your mouth open
\end{verbatim}

The \texttt{-\/-verbose\ (-v)} option causes \texttt{svn\ annotate} to
also include on each line the datestamp associated with that
line\textquotesingle s reported revision number. (This adds a
significant amount of width to each line of ouput, so
we\textquotesingle ll skip the demonstration here.)

As with \texttt{svn\ cat}, you can also ask \texttt{svn\ annotate} to
display previous versions of the file. This can be a handy trick when,
after finding out who most recently modified a particular line of
interest in the file, you then wish to see who modified the same line
prior to that.

\begin{verbatim}
$ svn blame rules.txt -r 2
     1      harry Be kind to others
     1      harry Freedom = Chocolate Ice Cream
     1      harry Everything in moderation
     1      harry Chew with your mouth open
\end{verbatim}

Unlike the \texttt{svn\ cat} command, the functionality of
\texttt{svn\ annotate} is tied heavily to the idea of ``lines'' of text
in a human-readable file. As such, if you attempt to run the command on
a file that Subversion has determined is \emph{not} human-readable (per
the file\textquotesingle s \texttt{svn:mime-type} property---see
\hyperref[svn.advanced.props.special.mime-type]{File Content Type} for
details), you\textquotesingle ll get an error message.

\begin{verbatim}
$ svn annotate images/logo.png
Skipping binary file (use --force to treat as text): 'images/logo.png'
$
\end{verbatim}

As revealed in the error message, you can use the \texttt{-\/-force}
option to disable this check and proceed with the annotation as if the
file\textquotesingle s contents are, in fact, human-readable and
line-based. Naturally, if you force Subversion to try to perform
line-based annotation on a nontextual file, you\textquotesingle ll get
what you asked for: a screenful of nonsense.

\begin{verbatim}
$ svn annotate images/logo.png --force
     6      harry \211PNG
     6      harry ^Z
     6      harry 
     7      harry \274\361\MI\300\365\353^X\300…
\end{verbatim}

Depending on your mood at the time you execute this command and your
reasons for doing so, you may find yourself typing
\texttt{svn\ blame\ …} or \texttt{svn\ praise\ …} instead of using the
canonical \texttt{svn\ annotate} command form. That\textquotesingle s
okay---the Subversion developers anticipated as much, so those
particular command aliases work, too!

Finally, as with many of Subversion\textquotesingle s informational
commands, you can also reference files in your \texttt{svn\ annotate}
command invocations by their repository URLs, allowing access to this
information even when you don\textquotesingle t have ready access to a
working copy.

\subsubsection{Listing versioned
directories}\label{svn.tour.history.browsing.list}

The \texttt{svn\ list} command shows you what files are in a repository
directory without actually downloading the files to your local machine:

\begin{verbatim}
$ svn list http://svn.example.com/repo/project
README
branches/
tags/
trunk/
\end{verbatim}

If you want a more detailed listing, pass the \texttt{-\/-verbose}
(\texttt{-v}) flag to get output like this:

\begin{verbatim}
$ svn list -v http://svn.example.com/repo/project
  23351 sally                 Feb 05 13:26 ./
  20620 harry            1084 Jul 13  2006 README
  23339 harry                 Feb 04 01:40 branches/
  23198 harry                 Jan 23 17:17 tags/
  23351 sally                 Feb 05 13:26 trunk/
\end{verbatim}

The columns tell you the revision at which the file or directory was
last modified, the user who modified it, the size if it is a file, the
date it was last modified, and the item\textquotesingle s name.

The \texttt{svn\ list} command with no arguments defaults to the
\emph{repository URL} of the current working directory, \emph{not} the
local working copy directory. After all, if you want a listing of your
local directory, you could use just plain \texttt{ls} (or any reasonable
non-Unixy equivalent).

\subsection{Fetching Older Repository
Snapshots}\label{svn.tour.history.snapshots}

In addition to all of the previous commands, you can use the
\texttt{-\/-revision} (\texttt{-r}) option with \texttt{svn\ update} to
take an entire working copy ``back in time'':\footnote{See? We told you
  that Subversion was a time machine.}

\begin{verbatim}
# Make the current directory look like it did in r1729.
$ svn update -r 1729
Updating '.':
…
$
\end{verbatim}

Many Subversion newcomers attempt to use the preceding
\texttt{svn\ update} example to ``undo'' committed changes, but this
won\textquotesingle t work as you can\textquotesingle t commit changes
that you obtain from backdating a working copy if the changed files have
newer revisions. See
\hyperref[svn.branchmerge.basicmerging.resurrect]{Resurrecting Deleted
Items} for a description of how to ``undo'' a commit.

If you\textquotesingle d prefer to create a whole new working copy from
an older snapshot, you can do so by modifying the typical
\texttt{svn\ checkout} command. As with \texttt{svn\ update}, you can
provide the \texttt{-\/-revision} (\texttt{-r}) option. But for reasons
that we cover in \hyperref[svn.advanced.pegrevs]{Peg and Operative
Revisions}, you might instead want to specify the target revision as
part of Subversion\textquotesingle s expanded URL syntax.

\begin{verbatim}
# Checkout the trunk from r1729.
$ svn checkout http://svn.example.com/svn/repo/trunk@1729 trunk-1729
…
# Checkout the current trunk as it looked in r1729.
$ svn checkout http://svn.example.com/svn/repo/trunk -r 1729 trunk-1729
…
$
\end{verbatim}

Lastly, if you\textquotesingle re building a release and wish to bundle
up your versioned files and directories, you can use
\texttt{svn\ export} to create a local copy of all or part of your
repository without any \texttt{.svn} administrative directories
included. The basic syntax of this subcommand is identical to that of
\texttt{svn\ checkout}:

\begin{verbatim}
# Export the trunk from the latest revision.
$ svn export http://svn.example.com/svn/repo/trunk trunk-export
…
# Export the trunk from r1729.
$ svn export http://svn.example.com/svn/repo/trunk@1729 trunk-1729
…
# Export the current trunk as it looked in r1729. 
$ svn export http://svn.example.com/svn/repo/trunk -r 1729 trunk-1729
…
$
\end{verbatim}

\section{Sometimes You Just Need to Clean Up}\label{svn.tour.cleanup}

Now that we\textquotesingle ve covered the day-to-day tasks that
you\textquotesingle ll frequently use Subversion for,
we\textquotesingle ll review a few administrative tasks relating to your
working copy.

\subsection{Disposing of a Working
Copy}\label{svn.tour.cleanup.disposal}

Subversion doesn\textquotesingle t track either the state or the
existence of working copies on the server, so there\textquotesingle s no
server overhead to keeping working copies around. Likewise,
there\textquotesingle s no need to let the server know that
you\textquotesingle re going to delete a working copy.

If you\textquotesingle re likely to use a working copy again,
there\textquotesingle s nothing wrong with just leaving it on disk until
you\textquotesingle re ready to use it again, at which point all it
takes is an \texttt{svn\ update} to bring it up to date and ready for
use.

However, if you\textquotesingle re definitely not going to use a working
copy again, you can safely delete the entire thing using whatever
directory removal capabilities your operating system offers. We
recommend that before you do so you run \texttt{svn\ status} and review
any files listed in its output that are prefixed with a \texttt{?} to
make certain that they\textquotesingle re not of importance.

\subsection{Recovering from an
Interruption}\label{svn.tour.cleanup.interruption}

When Subversion modifies your working copy---either your files or its
own administrative state---it tries to do so as safely as possible.
Before changing the working copy, Subversion logs its intentions in a
private ``to-do list'', of sorts. Next, it performs those actions to
effect the desired change, holding a lock on the relevant part of the
working copy while it works. This prevents other Subversion clients from
accessing the working copy mid-change. Finally, Subversion releases its
lock and cleans up its private to-do list. Architecturally, this is
similar to a journaled filesystem. If a Subversion operation is
interrupted (e.g, if the process is killed or if the machine crashes),
the private to-do list remains on disk. This allows Subversion to return
to that list later to complete any unfinished operations and return your
working copy to a consistent state.

This is exactly what \texttt{svn\ cleanup} does: it searches your
working copy and runs any leftover to-do items, removing working copy
locks as it completes those operations. If Subversion ever tells you
that some part of your working copy is ``locked,'' run
\texttt{svn\ cleanup} to remedy the problem. The \texttt{svn\ status}
command will inform you about administrative locks in the working copy,
too, by displaying an \texttt{L} next to those locked paths:

\begin{verbatim}
$ svn status
  L     somedir
M       somedir/foo.c
$ svn cleanup
$ svn status
M       somedir/foo.c
\end{verbatim}

Don\textquotesingle t confuse these working copy administrative locks
with the user-managed locks that Subversion users create when using the
lock-modify-unlock model of concurrent version control; see the sidebar
\hyperref[svn.advanced.locking.meanings]{The Many Meanings of ``Lock''}
for clarification.

\section{Dealing with Structural
Conflicts}\label{svn.tour.treeconflicts}

So far, we have only talked about conflicts at the level of file
content. When you and your collaborators make overlapping changes within
the same file, Subversion forces you to merge those changes before you
can commit.\footnote{Well, you \emph{could} mark files containing
  conflict markers as resolved and commit them, if you really wanted to.
  But this is rarely done in practice.}

But what happens if your collaborators move or delete a file that you
are still working on? Maybe there was a miscommunication, and one person
thinks the file should be deleted, while another person still wants to
commit changes to the file. Or maybe your collaborators did some
refactoring, renaming files and moving around directories in the
process. If you were still working on these files, those modifications
may need to be applied to the files at their new location. Such
conflicts manifest themselves at the directory tree structure level
rather than at the file content level, and are known as tree conflicts.

Tree conflicts prior to Subversion 1.6

Prior to Subversion 1.6, tree conflicts could yield rather unexpected
results. For example, if a file was locally modified, but had been
renamed in the repository, running \texttt{svn\ update} would make
Subversion carry out the following steps:

\begin{itemize}
\item
  Check the file to be renamed for local modifications.
\item
  Delete the file at its old location, and if it had local
  modifications, keep an on-disk copy of the file at the old location.
  This on-disk copy now appears as an unversioned file in the working
  copy.
\item
  Add the file, as it exists in the repository, at its new location.
\end{itemize}

When this situation arises, there is the possibility that the user makes
a commit without realizing that local modifications have been left in a
now-unversioned file in the working copy, and have not reached the
repository. This gets more and more likely (and tedious) if the number
of files affected by this problem is large.

Since Subversion 1.6, this and other similar situations are flagged as
conflicts in the working copy.

As with textual conflicts, tree conflicts prevent a commit from being
made from the conflicted state, giving the user the opportunity to
examine the state of the working copy for potential problems arising
from the tree conflict, and resolving any such problems before
committing.

\subsection{An Example Tree
Conflict}\label{svn.tour.treeconflicts.example}

Suppose a software project you were working on currently looked like
this:

\begin{verbatim}
$ svn list -Rv svn://svn.example.com/trunk/
     13 harry                 Sep 06 10:34 ./
     13 harry              27 Sep 06 10:34 COPYING
     13 harry              41 Sep 06 10:32 Makefile
     13 harry              53 Sep 06 10:34 README
     13 harry                 Sep 06 10:32 code/
     13 harry              54 Sep 06 10:32 code/bar.c
     13 harry             130 Sep 06 10:32 code/foo.c
$
\end{verbatim}

Later, in revision 14, your collaborator Harry renames the file
\texttt{bar.c} to \texttt{baz.c}. Unfortunately, you
don\textquotesingle t realize this yet. As it turns out, you are busy in
your working copy composing a different set of changes, some of which
also involve modifications to \texttt{bar.c}:

\begin{verbatim}
$ svn diff
Index: code/foo.c
===================================================================
--- code/foo.c  (revision 13)
+++ code/foo.c  (working copy)
@@ -3,5 +3,5 @@
 int main(int argc, char *argv[])
 {
     printf("I don't like being moved around!\n%s", bar());
-    return 0;
+    return 1;
 }
Index: code/bar.c
===================================================================
--- code/bar.c  (revision 13)
+++ code/bar.c  (working copy)
@@ -1,4 +1,4 @@
 const char *bar(void)
 {
-    return "Me neither!\n";
+    return "Well, I do like being moved around!\n";
 }
$
\end{verbatim}

You first realize that someone else has changed \texttt{bar.c} when your
own commit attempt fails:

\begin{verbatim}
$ svn commit -m "Small fixes"
Sending        code/bar.c
Transmitting file data .
svn: E155011: Commit failed (details follow):
svn: E155011: File '/home/svn/project/code/bar.c' is out of date
svn: E160013: File not found: transaction '14-e', path '/code/bar.c'
$
\end{verbatim}

At this point, you need to run \texttt{svn\ update}. Besides bringing
our working copy up to date so that you can see Harry\textquotesingle s
changes, this also flags a tree conflict so you have the opportunity to
evaluate and properly resolve it.

\begin{verbatim}
$ svn update
Updating '.':
   C code/bar.c
A    code/baz.c
U    Makefile
Updated to revision 14.
Summary of conflicts:
  Tree conflicts: 1
$
\end{verbatim}

In its output, \texttt{svn\ update} signifies tree conflicts using a
capital C in the fourth output column. \texttt{svn\ status} reveals
additional details of the conflict:

\begin{verbatim}
$ svn status
M       code/foo.c
A  +  C code/bar.c
      >   local edit, incoming delete upon update
Summary of conflicts:
  Tree conflicts: 1
$
\end{verbatim}

Note how \texttt{bar.c} is automatically scheduled for re-addition in
your working copy, which simplifies things in case you want to keep the
file.

Because a move in Subversion is implemented as a copy operation followed
by a delete operation, and these two operations cannot be easily related
to one another during an update, all Subversion can warn you about is an
incoming delete operation on a locally modified file. This delete
operation \emph{may} be part of a move, or it could be a genuine delete
operation. Determining exactly what semantic change was made to the
repository is important---you want to know just how your own edits fit
into the overall trajectory of the project. So read log messages, talk
to your collaborators, study the line-based differences---do whatever
you must do---to determine your best course of action.

In this case, Harry\textquotesingle s commit log message tells you what
you need to know.

\begin{verbatim}
$ svn log -r14 ^/trunk
------------------------------------------------------------------------
r14 | harry | 2011-09-06 10:38:17 -0400 (Tue, 06 Sep 2011) | 1 line
Changed paths:
   M /Makefile
   D /code/bar.c
   A /code/baz.c (from /code/bar.c:13)

Rename bar.c to baz.c, and adjust Makefile accordingly.
------------------------------------------------------------------------
$
\end{verbatim}

\texttt{svn\ info} shows the URLs of the items involved in the conflict.
The \emph{left} URL shows the source of the local side of the conflict,
while the \emph{right} URL shows the source of the incoming side of the
conflict. These URLs indicate where you should start searching the
repository\textquotesingle s history for the change which conflicts with
your local change.

\begin{verbatim}
$ svn info code/bar.c
Path: code/bar.c
Name: bar.c
URL: http://svn.example.com/svn/repo/trunk/code/bar.c
…
Tree conflict: local edit, incoming delete upon update
  Source  left: (file) ^/trunk/code/bar.c@4
  Source right: (none) ^/trunk/code/bar.c@5

$
\end{verbatim}

\texttt{bar.c} is now said to be the victim of a tree conflict. It
cannot be committed until the conflict is resolved:

\begin{verbatim}
$ svn commit -m "Small fixes" 
svn: E155015: Commit failed (details follow):
svn: E155015: Aborting commit: '/home/svn/project/code/bar.c' remains in confl
ict
$
\end{verbatim}

To resolve this conflict, you must either agree or disagree with the
move that Harry made.

If you agree with the move, your \texttt{bar.c} is superfluous.
You\textquotesingle ll want to delete it and mark the tree conflict as
resolved. But wait: you made changes to that file! Before deleting
\texttt{bar.c}, you need to decide if the changes you made to it need to
be applied elsewhere, for example to the new \texttt{baz.c} file where
all of \texttt{bar.c}\textquotesingle s code now lives.
Let\textquotesingle s assume that your changes do need to ``follow the
move''. Subversion isn\textquotesingle t smart enough to do this work
for you\footnote{In some cases, Subversion 1.5 and 1.6 \emph{would}
  actually handle this for you, but this somewhat hit-or-miss
  functionality was removed in Subversion 1.7.}, so you need to migrate
your changes manually.

In our example, you could manually re-make your change to \texttt{bar.c}
pretty easily---it was, after all, a single-line change.
That\textquotesingle s not always the case, though, so
we\textquotesingle ll show a more scalable approach.
We\textquotesingle ll first use \texttt{svn\ diff} to create a patch
file. Then we\textquotesingle ll edit the headers of that patch file to
point to the new name of our renamed file. Finally, we re-apply the
modified patch to our working copy.

\begin{verbatim}
$ svn diff code/bar.c > PATCHFILE
$ cat PATCHFILE
Index: code/bar.c
===================================================================
--- code/bar.c  (revision 14)
+++ code/bar.c  (working copy)
@@ -1,4 +1,4 @@
 const char *bar(void)
 {
-    return "Me neither!\n";
+    return "Well, I do like being moved around!\n";
 }
$ ### Edit PATCHFILE to refer to code/baz.c instead of code/bar.c
$ cat PATCHFILE
Index: code/baz.c
===================================================================
--- code/baz.c  (revision 14)
+++ code/baz.c  (working copy)
@@ -1,4 +1,4 @@
 const char *bar(void)
 {
-    return "Me neither!\n";
+    return "Well, I do like being moved around!\n";
 }
$ svn patch PATCHFILE
U         code/baz.c
$
\end{verbatim}

Now that the changes you originally made to \texttt{bar.c} have been
successfully reproduced in \texttt{baz.c}, you can delete \texttt{bar.c}
and resolve the conflict, instructing the resolution logic to accept
what is currently in the working copy as the desired result.

\begin{verbatim}
$ svn delete --force code/bar.c
D         code/bar.c
$ svn resolve --accept=working code/bar.c
Resolved conflicted state of 'code/bar.c'
$ svn status
M       code/foo.c
M       code/baz.c
$ svn diff
Index: code/foo.c
===================================================================
--- code/foo.c  (revision 14)
+++ code/foo.c  (working copy)
@@ -3,5 +3,5 @@
 int main(int argc, char *argv[])
 {
     printf("I don't like being moved around!\n%s", bar());
-    return 0;
+    return 1;
 }
Index: code/baz.c
===================================================================
--- code/baz.c  (revision 14)
+++ code/baz.c  (working copy)
@@ -1,4 +1,4 @@
 const char *bar(void)
 {
-    return "Me neither!\n";
+    return "Well, I do like being moved around!\n";
 }
$
\end{verbatim}

But what if you do not agree with the move? Well, in that case, you can
delete \texttt{baz.c} instead, after making sure any changes made to it
after it was renamed are either preserved or not worth keeping. (Do not
forget to also revert the changes Harry made to \texttt{Makefile}.)
Since \texttt{bar.c} is already scheduled for re-addition, there is
nothing else left to do, and the conflict can be marked resolved:

\begin{verbatim}
$ svn delete --force code/baz.c
D         code/baz.c
$ svn resolve --accept=working code/bar.c
Resolved conflicted state of 'code/bar.c'
$ svn status
M       code/foo.c
A  +    code/bar.c
D       code/baz.c
M       Makefile
$ svn diff
Index: code/foo.c
===================================================================
--- code/foo.c  (revision 14)
+++ code/foo.c  (working copy)
@@ -3,5 +3,5 @@
 int main(int argc, char *argv[])
 {
     printf("I don't like being moved around!\n%s", bar());
-    return 0;
+    return 1;
 }
Index: code/bar.c
===================================================================
--- code/bar.c  (revision 14)
+++ code/bar.c  (working copy)
@@ -1,4 +1,4 @@
 const char *bar(void)
 {
-    return "Me neither!\n";
+    return "Well, I do like being moved around!\n";
 }
Index: code/baz.c
===================================================================
--- code/baz.c  (revision 14)
+++ code/baz.c  (working copy)
@@ -1,4 +0,0 @@
-const char *bar(void)
-{
-    return "Me neither!\n";
-}
Index: Makefile
===================================================================
--- Makefile    (revision 14)
+++ Makefile    (working copy)
@@ -1,2 +1,2 @@
 foo: 
-   $(CC) -o $@ code/foo.c code/baz.c
+   $(CC) -o $@ code/foo.c code/bar.c
\end{verbatim}

You\textquotesingle ve now resolved your first tree conflict! You can
commit your changes and tell Harry during tea break about all the extra
work he caused for you.

\section{Summary}\label{svn.tour.summary}

Now we\textquotesingle ve covered most of the Subversion client
commands. Notable exceptions are those dealing with branching and
merging (see \hyperref[svn.branchmerge]{Branching and Merging}) and
properties (see \hyperref[svn.advanced.props]{Properties}). However, you
may want to take a moment to skim through \hyperref[svn.ref.svn]{svn
Reference---Subversion Command-Line Client} to get an idea of all the
different commands that Subversion has---and how you can use them to
make your work easier.

